\clearpage
\phantomsection
\addcontentsline{toc}{chapter}{Cronograma de 15 Encontros e Plano de Estudos}
\chapter*{Cronograma de 15 Encontros e Plano de Estudos}
\markboth{CRONOGRAMA — 15 ENCONTROS}{CRONOGRAMA — 15 ENCONTROS}

O componente é desenvolvido em **15 encontros presenciais de 4 horas** (total de 60 h). O quadro a seguir integra cada aula do curso ao respectivo capítulo desta apostila e às entregas avaliativas (AV1, AV2 e AV3).

\renewcommand{\arraystretch}{1.22}
\small
\begin{longtable}{@{}c p{2.8cm} p{5.7cm} p{4.3cm}@{}}
  \toprule
  \rowcolor{etIndigo}
  \textcolor{white}{\bfseries Encontro} & \textcolor{white}{\bfseries Unidade / Cap.} & \textcolor{white}{\bfseries Tema Central da Aula} & \textcolor{white}{\bfseries Autores-Chave e Avaliação} \\
  \midrule
  \endfirsthead
  \toprule
  \rowcolor{etIndigo}
  \textcolor{white}{\bfseries Encontro} & \textcolor{white}{\bfseries Unidade / Cap.} & \textcolor{white}{\bfseries Tema Central da Aula} & \textcolor{white}{\bfseries Autores-Chave e Avaliação} \\
  \midrule
  \endhead
  \bottomrule
  \endfoot

  \textbf{1} & \textbf{Unid. I — Cap. 1} & Empreendedorismo de Base Tecnológica e Desenvolvimento Territorial & Schumpeter; Kuhn; Dosi; Milton Santos. \\
  \rowcolor{bgGray}
  \textbf{2} & \textbf{Unid. I — Cap. 2} & Estruturas de Mercado, Barreiras de Entrada e Aglomerações Produtivas & Porter; Eficiência Coletiva (APLs). \\
  \textbf{3} & \textbf{Unid. I — Cap. 3} & Ecossistemas Empreendedores e Sistemas Regionais de Inovação & Isenberg; Stam; Boschma (Proximidades). \\
  \rowcolor{bgGray}
  \textbf{4} & \textbf{Unid. I — Cap. 4} & Governança Territorial da Inovação e Políticas Públicas de CT\&I & Etzkowitz/Leydesdorff; Mazzucato. \\
  \textbf{5} & \textbf{Unid. I — Cap. 5} & Empreendedorismo Social e Negócios de Impacto Territorial & Austin et al. (PCDO); Teoria da Mudança. \\
  \rowcolor{bgTeal}
  \textbf{6} & \textbf{Unid. I — Cap. 6} & Inovação Social, Economia Solidária, Bioeconomia e IGs & Dagnino (Tec. Social); INPI. \textbf{\color{etAlert}AV1 (30\%) + Simulado I.} \\
  \midrule
  \textbf{7} & \textbf{Unid. II — Cap. 7} & Processo Empreendedor e Criação de Novos Negócios & Shane \& Venkataraman; Sarasvathy. \\
  \rowcolor{bgGray}
  \textbf{8} & \textbf{Unid. II — Cap. 8} & Identificação e Análise de Oportunidades Tecnológicas e Territoriais & Alvarez \& Barney; Drucker; TAM/SAM/SOM; PI. \\
  \textbf{9} & \textbf{Unid. II — Cap. 9} & Modelos de Negócio, Proposta de Valor e Ciclo Lean & Chesbrough \& Rosenbloom; Osterwalder; Blank/Ries. \\
  \rowcolor{bgGray}
  \textbf{10} & \textbf{Unid. II — Cap. 10} & Planos de Negócio e Instrumentos de Fomento Regional & Escada do Capital; FINEP/BNDES/BNB/FAPESB. \\
  \textbf{11} & \textbf{Unid. II — Cap. 11} & Empreendedorismo Corporativo, Intraempreendedorismo e Spin-offs & Pinchot; Guth \& Ginsberg; Chesbrough. \textbf{Simulado II.} \\
  \midrule
  \textbf{12} & \textbf{Unid. III — Cap. 12} & Elaboração de Projetos de Investimento e Decisão Locacional & Woiler \& Mathias; CAPEX/OPEX; Matriz Locacional. \\
  \rowcolor{bgGray}
  \textbf{13} & \textbf{Unid. III — Cap. 13} & Métodos e Ferramentas para Análise de Investimentos & VPL, TIR, MTIR, VAUE, Payback, Sensibilidade. \\
  \textbf{14} & \textbf{Unid. III — Cap. 14} & Valuation de Empreendimentos Tecnológicos e Projetos de Impacto & Damodaran (FCD); VC; Opções Reais. \textbf{\color{etAlert}AV2 (20\%) + Simulado III.} \\
  \midrule
  \rowcolor{bgAmber}
  \textbf{15} & \textbf{Unid. IV — Cap. 15} & Seminário Final de Integração: Ecossistemas, Negócios e Viabilidade Territorial & Defesa Oral + Rubrica \textbf{\color{etAlert}AV3 (50\%)} + Avaliação Final. \\
\end{longtable}
\normalsize
